\subsection{Universal set of quantum gates}

After studying Clifford and non-Clifford gates, the next question is: \emph{what is the minimum set of gates needed to build \textbf{any quantum computation}?} This leads to the concept of a \textbf{universal set of quantum gates}, which is a collection of gates that can be combined to approximate any unitary operation on a quantum system to arbitrary precision.

\highspace
\begin{definitionbox}[: Universal Set of Quantum Gates]
    A \textbf{set of quantum gates} is called \definitionWithSpecificIndex{universal}{Universal}{Universal Set of Quantum Gates} if any \hl{unitary transformation can be approximated to arbitrary accuracy using only gates from that set.}

    \highspace
    In other words, a set of gates is universal if using only gates from that set, we can approximate any unitary transformation to arbitrary precision.
    \begin{equation*}
        \text{Universal gate set} \Longrightarrow \forall U,\ \exists\ \text{a circuit approximating } U
    \end{equation*}
\end{definitionbox}

\begin{flushleft}
    \textcolor{Green3}{\faIcon{question-circle} \textbf{Why ``Approximate'' and not ``Exactly''?}}
\end{flushleft}
In quantum computing, every quantum gate is represented by a \textbf{unitary matrix}. For example the Hadamard gate, the $T$ gate, and the CNOT gate. However, the set of all unitary matrices is \textbf{continuous} (i.e., there are infinitely many unitary transformations). For example, a single-qubit rotation may involve any real angle $R_z(\theta)$ with infinitely many possible values of $\theta$:
\begin{equation*}
    R_z(\theta) = \begin{bmatrix}
        e^{-i\theta/2} & 0 \\
        0 & e^{i\theta/2}
    \end{bmatrix}
\end{equation*}
Here, the angle $\theta$ can be $0.1$, $0.12345$, $\pi/7$ or any other real number. There are infinitely many real numbers, so there are infinitely many possible rotations and therefore infinitely many possible unitary matrices.

\highspace
\textcolor{Green3}{\faIcon{question-circle} \textbf{Why not build hardware for every matrix?}} Building a quantum gate for every possible unitary transformation is \textbf{impossible}. Imagine trying to crate one physical instruction for every possible angle:
\begin{itemize}
    \item One gate for $0.1$ radians
    \item One gate for $0.12345$ radians
    \item One gate for $\pi/7$ radians
    \item One gate for $0.00001$ radians
\end{itemize}
And so on forever. This is clearly impractical.

\highspace
\textcolor{Green3}{\faIcon{check-circle} \textbf{The solution: a Finite Gate Set.}} Instead of trying to build a gate for every possible unitary, we choose a \textbf{small set of standard gates}, for example $\left\{H, \, \text{CNOT}, \, T\right\}$. These are the only \textbf{primitive gates the machine needs to implement directly}. All other operations are constructed by combining them.

\newpage

\begin{flushleft}
    \textcolor{Red2}{\faIcon{exclamation-triangle} \textbf{Which gates are in the universal set?}}
\end{flushleft}
To keep things simple, we could choose Clifford gates, such as $H$ and CNOT, as our set of gates. Although these gates are powerful, they can still be efficiently simulated on a classical computer (\autopageref{question:why-non-clifford-gates-important}). Therefore, \hl{Clifford gates alone are not sufficient for universal quantum computation.} For this reason, in general, a \textbf{universal set requires at least one non-Clifford gate}, and the most common choice is the $T$ gate because of its simplicity and its relationship to the $S$ gate (\autopageref{sec:non-clifford-gates}). A \hl{widely used universal set of quantum gates is:}
\begin{equation}
    \left\{H, \, \text{CNOT}, \, T\right\}
\end{equation}
\textcolor{Green3}{\faIcon{question} \textbf{Why these three gates are enough?}}
\begin{itemize}
    \item \textbf{Hadamard} $H$ \hl{transforms basis states into superpositions.} For example, the \ket{0} state becomes the plus state (\autopageref{eq:plus-state}):
    \begin{equation*}
        H\ket{0} = \ket{+} = \frac{1}{\sqrt{2}}(\ket{0} + \ket{1})
    \end{equation*}

    \item \textbf{CNOT} creates \hl{entanglement between qubits.} For example, applying the CNOT to the state $\ket{+0}$ creates the Bell state $\ket{\Phi^+}$:
    \begin{equation*}
        \text{CNOT}(\ket{+0}) = \ket{\Phi^+} = \frac{1}{\sqrt{2}}(\ket{00} + \ket{11})
    \end{equation*}

    \item \textbf{$T$ gate} adds a \hl{$45\degree$ phase and breaks the Clifford-only restriction}, allowing us to access a much larger set of unitary transformations. Because $T$ is a non-Clifford gate, it enables us to perform operations that cannot be efficiently simulated classically.
\end{itemize}
Hadamard changes basis and creates superposition, CNOT creates entanglement between qubits, and $T$ provides the non-Clifford component required for universality. Together, \textbf{they can approximate any unitary operation}.

\highspace
Other universal gate sets exist, such as:
\begin{itemize}
    \item Toffoli and Hadamard: $\left\{H, \, \text{Toffoli}\right\}$.
    \item Single-qubit rotations and CNOT: $\left\{R_x(\theta), \, R_y(\theta), \, R_z(\theta), \, \text{CNOT}\right\}$.
\end{itemize}

\highspace
\begin{flushleft}
    \textcolor{Green3}{\faIcon{book} \textbf{How to Generate a Universal Set of Quantum Gates}}
\end{flushleft}
A universal gate set is obtained by combining two ingredients:
\begin{enumerate}
    \item \hl{A set of gates that can generate arbitrary single-qubit unitary matrices.}
    \item \hl{One entangling two-qubit gates}, usually the CNOT gate (\autopageref{sec:controlled-not-cnot-gate}), which can create entanglement between qubits.
\end{enumerate}
If both conditions are satisfied, the gate set is universal. Let's see why this is the case:
\begin{enumerate}
    \item \textbf{Choose Gates that Generate Arbitrary Single-Qubit Unitaries} (e.g., $H$ and $T$ gates). A single qubit can undergo any unitary transformation:
    \begin{equation*}
        U \in U^{2} = \left\{U \in \mathbb{C}^{2 \times 2} : U^\dagger U = I\right\}
    \end{equation*}
    To obtain universality, we need to approximate any such $2 \times 2$ unitary matrix. A common way is to use a Hadamard gate and a $T$ gate. The pair of gates $\{H, T\}$ can approximate any single-qubit unitary to arbitrary precision.

    \textcolor{Green3}{\faIcon{question-circle} \textbf{Why universality?}} While mathematically there are infinitely many possible $2 \times 2$ unitary matrices, our hardware can implement only a \textbf{finite, discrete set} of gates exactly (similar to how a digital computer uses a finite instruction set). Therefore, the goal of a \emph{universal} gate set is not to implement every unitary matrix exactly. Instead, \hl{for any target unitary matrix $U$ and for any desired accuracy $\varepsilon > 0$, there must exist a finite sequence of gates from the set whose overall operation $V$ satisfies:}\footnote{%
        The matrix $U-V$ contains the element-by-element differences, but by itself it is still a matrix, not a single number. To obtain one number representing the size of the difference, we apply a \textbf{norm}. It plays exactly the same role as absolute value for numbers. In other words, it means that the implemented operation $V$ is within distance $\varepsilon$ of the desired operation $U$.
    }
    \begin{equation*}
        \left\|U - V\right\| < \varepsilon
    \end{equation*}
    In other words, we can \hl{make the approximation error as small as we want}. This ability to approximate any unitary operation to arbitrary precision is what makes a gate set universal.


    \item \textbf{Add an Entangling (two-qubit) Gate}. Single-qubit gates alone cannot create entanglement (\autopageref{question:why-single-qubit-gates-cannot-create-entanglement}). To manipulate multiple qubits, we \hl{add an entangling gate such as CNOT}. Once CNOT is available, qubits can interact and become entangled.
\end{enumerate}
Combining these two ingredients, the resulting set is universal, meaning that any quantum computation can be approximated using just these gates.